\chapter{Transcripciones de las entrevistas}
\label{anexo:entrevistas}

Criterio de depuración. Las siguientes transcripciones se presentan en formato de diálogo directo, en primera persona para entrevistadores y profesionales entrevistados. Se eliminaron pruebas técnicas de cámara o audio, repeticiones sin valor analítico, muletillas excesivas y errores evidentes de la transcripción automática. No se reemplazaron las respuestas por síntesis en tercera persona, ya que el anexo funciona como registro primario del relevamiento cualitativo. Cuando una frase del audio resultaba dudosa, se conservó el sentido más probable sin agregar conclusiones externas.

Identificación de participantes. Los nombres e instituciones se consignan según la información provista por los autores del PFI para esta versión de trabajo. El cuerpo de cada entrevista mantiene el intercambio en primera persona y no incluye análisis interpretativo; dicho análisis se desarrollará en el capítulo de user research. % CROSS_REFERENCE_MAPPING_REQUIRED

% SOURCE_NUMBERING_ANOMALY_PRESERVED: E.1 literal in DOCX Heading2
\phantomsection
\section*{E.1 Entrevista 1 (EN-01) - Dra. Susana Torres (M.N. 179537)}
\addcontentsline{toc}{section}{E.1 Entrevista 1 (EN-01) - Dra. Susana Torres (M.N. 179537)}
\label{anexo:entrevista-en01}

\noindent\textbf{Dra. Susana Torres:} Ustedes serían la parte de programación específica, o sea, desarrollar la idea. ¿Qué es lo que necesitan? ¿Ciertos parámetros? ¿Qué tienen que tener en cuenta al momento de evaluar una persona en columna lumbar? Más que nada, si es para estenosis, cuántas patologías abarcarían. Me pregunté un montón de cosas.

\noindent\textbf{Enzo Asplanatti:} Por ahora nosotros mantenemos la parte ingenieril, que sería automatizar ese proceso y, con ayuda de la computadora, lograr encontrar la mayor cantidad de hallazgos posibles. Todavía no tenemos el alcance completamente acotado a patologías específicas. Por ahora lo estamos achicando a esto.

\noindent\textbf{Francisco Fabrello:} Estamos esperando la respuesta de la cátedra. Presentamos el proyecto, nos dijeron que avancemos y que después nos iban a dar correcciones finas. Entonces, por ahora, la idea es entrenar el modelo con datos ya etiquetados para que aprenda a detectar, y después llevarlo hacia lo que nos recomienden. También nos pidieron validar con profesionales, haciendo entrevistas con algunas preguntas, para ver qué opinan y qué podríamos agregar o quitar.

\noindent\textbf{Dra. Susana Torres:} ¿Ustedes tienen una lista de preguntas?

\noindent\textbf{Francisco Fabrello:} Sí, más o menos. La primera sería: ¿cómo trabajás con resonancias y qué lugar ocupan dentro de tu práctica?

\noindent\textbf{Dra. Susana Torres:} Resonancias lumbares, hoy por hoy estoy trabajando en FLENI. El principal volumen son los cerebros, pero después le sigue seguro la resonancia lumbar por hernias discales. Hay otros hallazgos no pensados dentro del diagnóstico presuntivo, pero la mayoría es por radiculopatía o compresión radicular.

\noindent\textbf{Francisco Fabrello:} Entonces, ¿cuál sería más o menos el caso que encontrás más frecuente?

\noindent\textbf{Dra. Susana Torres:} Protrusiones discales.

\noindent\textbf{Francisco Fabrello:} ¿Y eso sería, en palabras más simples?

\noindent\textbf{Dra. Susana Torres:} La hernia discal es como la entidad general. Hay distintos grados: está el abombamiento, la protrusión y la extrusión. La protrusión es la más frecuente.

\noindent\textbf{Francisco Fabrello:} Perfecto. Y cuando tenés que hacer el análisis de una resonancia, ¿cómo sería el recorrido para visualizarla?

\noindent\textbf{Dra. Susana Torres:} Lo que ya viene estudiado, que también hace poco se usa en FLENI, es más que nada la parte de volumetría encefálica, que tira resultados de cuantificación de sustancia gris, sustancia blanca, etcétera; pero para resonancia lumbar todavía no he trabajado con eso. Mi orden de trabajo generalmente es dividir la pantalla. El axial T2 es el que siempre estoy corroborando y en la otra pantalla veo el sagital T2, el sagital T1 y el sagital STIR para ver si hay edema óseo o cambios degenerativos a nivel de los cuerpos vertebrales. Una vez visto eso, los espacios y la alineación, paso al sagital T2 y al axial T2. En el axial puedo ver si hay alguna compresión, cómo está el canal, si hay alguna estrechez a nivel del canal y cómo están los neuroforámenes.

\noindent\textbf{Francisco Fabrello:} Perfecto. ¿Qué considerás más importante para que el informe sea útil?

\noindent\textbf{Dra. Susana Torres:} Tener en cuenta la clínica del paciente y el diagnóstico presuntivo. Eso es lo que le está doliendo. Puede haber un montón de cambios degenerativos, pero si el paciente viene por una hernia L5-S1, voy enfocada ahí. Ahí hago mi conclusión: dónde le duele y por qué le duele.

\noindent\textbf{Francisco Fabrello:} Si este sistema pudiera delimitar automáticamente las vértebras, los discos, el canal espinal y, a partir de eso, calcular mediciones como la altura de discos, el diámetro AP del canal o la alineación, ¿qué utilidad real podría tener o qué dudas te generaría?

\noindent\textbf{Dra. Susana Torres:} A eso que nombraste le agregaría, además de la altura de los discos, la intensidad de señal de los discos. Pueden tener altura respetada, pero ya verse signos de deshidratación, que es súper común y así empiezan los cambios degenerativos. El diámetro AP del canal, en cuanto a lumbar, sería más que nada el área. Muchas veces hay cambios degenerativos en los elementos posteriores y facetas que se hipertrofian. Eso genera una disminución del diámetro transverso del canal. Entonces, por ahí el AP está bien, pero como tenés cambios degenerativos atrás, las raíces de la cola de caballo se apiñan y eso ya es una estenosis leve. También hay que tener en cuenta que el paciente está acostado; con ciertos movimientos o parado puede cambiar. En cuanto a la alineación, me parece re importante, porque el paciente está acostado y no tiene la sobrecarga de la gravedad. La alineación posterior también es súper importante para sospechar algún tipo de lisis o listesis.

\noindent\textbf{Francisco Fabrello:} Si el sistema trabajara principalmente con cortes sagitales T1 y T2, pero no con cortes axiales, ¿qué podría evaluarse y qué quedaría incompleto o riesgoso?

\noindent\textbf{Dra. Susana Torres:} Bastante quedaría incompleto. Al hacer el análisis nivel por nivel, si me hacés elegir, es el sagital T2 y el axial T2. Ahí veo también la extensión hacia los neuroforámenes. Incluir la estenosis a nivel de los neuroforámenes es re importante y se ve mejor correlacionando los dos planos.

\noindent\textbf{Francisco Fabrello:} ¿Qué errores de una herramienta de este tipo serían graves o inaceptables desde el punto de vista clínico?

\noindent\textbf{Dra. Susana Torres:} Los mismos que tenemos nosotros, básicamente: sobredimensionar algo o infravalorar algún hallazgo. Lo importante está todo en la calidad de los datos que le carguen, básicamente. Y qué ground truth es la que utilizan.

\noindent\textbf{Francisco Fabrello:} Claro. Yendo más al lado de la interfaz, ¿qué tendría que mostrar o permitir el sistema para que resulte confiable y no agregue ruido o haga más tediosa la rutina?

\noindent\textbf{Dra. Susana Torres:} Pienso en lo que evalúo a diario. Por ejemplo, si me dice que tiene una estenosis grado 2 de los neuroforámenes, poder verlo. Que me lo marque primero y, si estoy de acuerdo, lo dejo; si no, hago una corrección manual y digo que es grado 1 o grado 2. Igual ahí también hay sesgo, porque si una IA te dice que es grado 2 y quizás es grado 1, depende de la experiencia y de cómo se interprete el software.

\noindent\textbf{Francisco Fabrello:} Para que este software sea viable desde el punto de vista clínico, ¿contra qué tendríamos que validarlo? Por ejemplo, segmentaciones manuales de especialistas o mediciones de expertos.

\noindent\textbf{Dra. Susana Torres:} Sí, segmentación manual de especialistas sería ideal. Yo trabajo para una empresa que se llama Prenuvo, que hace total body segmentations con resonancias y distintos tipos de informes para hallazgos positivos. Las segmentaciones las hacemos entre médicos, pero siempre hay como un pool; cada estudio está tres veces chequeado. La calidad de los datos me parece lo más importante. Se nota cuando el que primero hizo la revisión es un médico general, un residente de radiología o un especialista.

\noindent\textbf{Francisco Fabrello:} ¿Qué alcance te parecería más defendible? ¿Qué sacarías para no prometer de más?

\noindent\textbf{Dra. Susana Torres:} Lo acotaría. Como les preguntaba al principio: ¿cuál era el objetivo? ¿Valorar estenosis? ¿Delimitar estenosis de canal y estenosis neuroforaminal? Porque si no hay un montón de patologías y un montón de cuestiones que se pueden evaluar. Generalmente, si me preguntás qué es lo que más vemos y lo que generalmente causa dolor o consulta del paciente, son cambios degenerativos. También cambios tipo Modic de las plataformas. Lo acotaría, porque si no va a ser muy difícil abarcar todo y se les va a hacer eterno.

\noindent\textbf{Francisco Fabrello:} ¿Alguna red flag que te haría decir que el proyecto no debería plantearse así?

\noindent\textbf{Dra. Susana Torres:} No veo que tuviera alguna red flag. La verdad es un tema bastante estudiado. Ya hay distintos modelos e instituciones afuera que están trabajando así, con informes y resultados estructurados. Ayudaría un montón. Recalco nuevamente la calidad de los datos de entrenamiento, pero el resto me parece fantástico.

\noindent\textbf{Francisco Fabrello:} ¿Qué dirías que sería un diferencial del proyecto? Por ejemplo, que genere el informe y muestre paso a paso lo que hizo, o que deje visualizar la resonancia con el análisis encima.

\noindent\textbf{Dra. Susana Torres:} Visualmente, creo que hablo por la mayoría de los radiólogos: nosotros nos guiamos mucho por lo visual. Si las segmentaciones están como máscaras dentro de los hallazgos y yo veo que tiene pérdida de altura del disco y me muestran dónde lo está viendo, eso me genera confianza. Lo mismo con las estenosis, ya sea del canal o de los forámenes.

\noindent\textbf{Francisco Fabrello:} ¿Qué opinás de que el sistema permita cargar una resonancia anterior y una actual para hacer una comparativa?

\noindent\textbf{Dra. Susana Torres:} Buenísimo. Eso ayudaría a decir qué era lo que le dolía antes, qué aumentó y cómo fueron evolucionando esos cambios degenerativos. Sumaría un montón.

\noindent\textbf{Enzo Asplanatti:} ¿Cómo continúa después, si avanza el proyecto?

\noindent\textbf{Dra. Susana Torres:} Si el proyecto continúa y progresa, yo me dedico a segmentación y control de calidad, así que tienen mi contacto. Me interesa mucho la inteligencia artificial porque es lo que se viene. Cualquier cosa que necesiten, tienen mi contacto.

\noindent\textbf{Francisco Fabrello:} La idea es que primero validemos con profesionales la viabilidad y después, más adelante, con un producto semiarmado, podamos volver a hacer preguntas tipo checklist del uno al cinco para evaluar qué mejorar.

\noindent\textbf{Dra. Susana Torres:} Claro, cuál es la practicidad en el día a día.

\noindent\textbf{Francisco Fabrello:} También pensamos que el sistema muestre la máscara encima, una tabla de datos y que el profesional pueda corregir el dato que predijo, porque vimos que el modelo a veces falla.

\noindent\textbf{Dra. Susana Torres:} O sea, que pasa.

\noindent\textbf{Francisco Fabrello:} Por ejemplo, mide el canal, lo pinta, pero le falta un pedacito. Entiendo que depende de la resonancia, de si es T1 o T2, y de qué estructuras se ven mejor.

\noindent\textbf{Dra. Susana Torres:} Claro. En cuanto a lo lumbar, les recalco que no solamente importa el diámetro, sino también la pérdida de líquido entre las raíces. A nivel cervical puede ser solamente diámetro, pero a nivel lumbar tenés la cola de caballo. Entonces hay que medir cuán juntas están las raíces y si se sigue viendo líquido entre ellas.

\noindent\textbf{Francisco Fabrello:} Lo vamos a tener en cuenta para entrenar la IA.

\noindent\textbf{Enzo Asplanatti:} Eso está bueno.

\noindent\textbf{Dra. Susana Torres:} Bueno, muchas gracias.

\noindent\textbf{Francisco Fabrello:} Gracias. ¿En un futuro te podríamos escribir para seguir validando?

\noindent\textbf{Dra. Susana Torres:} Sí, encantada, ningún problema. Me interesa mucho.

\noindent\textbf{Francisco Fabrello:} Muchas gracias.

\noindent\textbf{Dra. Susana Torres:} Dale, nos vemos.

\noindent\textbf{Enzo Asplanatti:} Chau, chau.

% MIGRATION_PENDING: E.2 / EN-02
% MIGRATION_PENDING: E.3 / EN-03
% MIGRATION_PENDING: E.4 / EN-04
% MIGRATION_PENDING: E.5 / EN-05
